% PASS20260921 physical FI factor routing through E6+A2.
\subsection*{The FI three-block is the external $A_2$; the physical $SU(5)$ lies in $E_6$}

Combining the exact physical twin-$A_4$ embedding with the FI-induced
$E_6+A_2$ grading gives an objectwise root-set routing.  Of the eight
FI-neutral roots inside the hypercharge-orthogonal organizer $A_4$, the six
roots belonging to its three-block are exactly the six-root external $A_2$
component.  The remaining two roots form an $A_1$ lying inside the $72$-root
$E_6$ component.  All $20$ roots of the physical Wilson-line $SU(5)$ $A_4$
also lie in that same $E_6$ component.

Hence, at root-system level,
\[
A_4^{\rm gauge}+A_1^{\rm FI}\subset E_6,
\qquad
A_2^{\rm FI}=A_2^{\rm external},
\]
with the two internal factors orthogonal and of combined rank five.  This is
compatible with the standard $SU(5)\times SU(2)\times U(1)\subset E_6$
chain, while the FI three-block is the external $SU(3)$ factor.

\paragraph{Scope firewall.}
The repository also contains separately constructed W33/affine
$SU(3)_{\rm family}$ actions.  The theorem above identifies an exact physical
$E_8$ root $A_2$, but does not identify it with those canonical family actions
without an explicit Weyl/intertwiner map.

The executable witness is
\texttt{analysis/w33\_physical\_fi\_factor\_routing\_e6\_a2.py}.
